\begin{theorem}[Connected complement of a contiguous simply connected block]%
    \label{thm:peps_torus_simply_connected_complement}
    \lean{TNLean.PEPS.torusGraph_compl_connected_of_isSimplyConnected,
        TNLean.PEPS.nonempty_torusRegionBoundaryEdge_of_isSimplyConnected}
    \leanok
    Let both torus circumferences be at least three. Let $R$ be a finite
    set of sites whose induced nearest-neighbour graph is connected and
    whose actual closed-cell realization is simply connected. Then the
    induced graph on the complementary sites is connected, and at least
    one bond crosses between $R$ and its complement.

    Consequently, spanning trees and roots on both sides of the cut,
    and a numbering of the crossing bonds by $n+1$ labels, may be chosen
    without further hypotheses. Occupied graph connectivity is distinct
    from simple connectedness of the closed-cell union: cells meeting
    only at a corner need not be connected by lattice bonds. This is the
    contiguous-block geometry used in Section~6.3 of
    \cite{Schuch2010PEPS}, rather than a claim that every simply connected
    closed-cell subset is a contiguous block.
\end{theorem}

\begin{proof}\leanok
    \uses{thm:peps_torus_integer_cell_lift,
        thm:peps_integer_exterior_collar_paths,
        thm:peps_integer_boundary_contour_connected,
        thm:peps_torus_complement_path_replacement,
        thm:peps_torus_rectangle_cut_connectivity}
    Lift the occupied region to a simply connected finite integer-cell
    union whose exterior collar projects outside the region. The exposed
    contour joins its nearby missing centres within the safe exterior
    graph. These joins yield complementary lattice walks between all
    outside endpoints of crossing bonds. Every complementary site has
    a walk to an occupied site in the full torus graph; its first crossing
    joins it, within the complement, to such an endpoint. Hence all
    complementary sites are connected. The collar also supplies a
    complementary site, so a walk from an occupied site crosses at least
    one boundary bond.
\end{proof}
